\subsubsection{Фильтр Блюма и подобные}

Фильтр Блюма - вероятностный алгоритм, для проверки принадлежности элемента множеству, гарантирующий отсутствие ложно-отрицательных срабатываний.

Устройство фильтра: есть фиксированное $m$ - количество битов и $k$ и $k$ хэш функций $h_1, \ldots, h_k $, из некоторого семейства, $h_i:\mathcal{K}\rightarrow \{0, \ldots, m-1 \}$. Положим массив $A$ битов размера $m$. Когда в множество добавляется новый ключ $K$, положим $\forall i \in \{1, \ldots, k\} A[h_i(K)] := 1$. Проверка принадлежности ключа множеству проводится следующим образом: необходимо проверить, что $\forall i A[h_i(K)] == 1$.


Фильтр Блюма позволяет проверить отсутствие ключа в SSTable, перед началом дорого (с точки зрения числа I/O запросов) поиска ключа в самой таблице. В свою очередь, ложно-положительные срабатывания не влияют на корректность ответа, только лишь добавляют расходы на поиск внутри самой компоненты. Вероятность такого  ложно-положительного срабатывания фильтра не превосходит $(1 - e ^ {-kn / m})^k$, тогда оптимальное количество хэш  функций для минимизации ошибки достигается при $k:=\frac{m}{2}\ln 2$, что на практике оказывается ~10 и позволяет достичь $1\%$ ложно-положительных срабатываний. А так как фильтр Блюма очень маленький и может быть загружен в оперативную память, накладные расходы на $I/O$ запросы к диску во многих сценариях получается существенно сократить.

Собственно, идея о том, что с помощью некоторой небольшой вероятностной структуры можно с хорошей точностью быстро проверить отсутствие ключа (если ключа нет в таблице, то структура должна гарантированно это сообщить) в одной компоненте, находит применение во многих современных СУБД. LevelDB использует встроенный фильтр Блюма - NewBloomFilterPolicy \cite{leveldbgit}. А RocksDB позволяет выбирать фильтр: NewBloomFilterPolicy, NewRibbonFilterPolicy (Ribbon Filter - некоторый аналог фильтра блюма, который занимает меньше памяти).
 \cite{rocksdbgit}.